\unit{Exposé 195 -- Descent Technique II}
\sid{195}
\src{369}
\seminaren{12th year, 1959/60, no. 195\\February 1960}
\paperhead{DESCENT TECHNIQUE AND EXISTENCE THEOREMS IN ALGEBRAIC GEOMETRY\\
II. THE EXISTENCE THEOREM IN FORMAL MODULI THEORY}%
  {by Alexander GROTHENDIECK}

\fsection{A}{Representable and pro-representable functors}

\subsection*{1. Representable functors}
\addcontentsline{toc}{subsection}{1. Representable functors}

Let $\mathcal C$ be a category. For every $X\in\mathcal C$, let $h_X$ be the
contravariant functor from $\mathcal C$ to the category
$(\operatorname{Ens})$ of sets,
\[
  h_X:\mathcal C\longrightarrow(\operatorname{Ens}),
\]
defined by the formula
\[
  h_X(Y)=\operatorname{Hom}(Y,X).
\]
A morphism $X\longrightarrow X'$ in $\mathcal C$ gives, in the evident way, a
homomorphism of functors $h_X\longrightarrow h_{X'}$: $h_X$ is a covariant
functor of $X$. In other words, we have defined a canonical covariant functor
\[
  h:\mathcal C\longrightarrow
  \operatorname{Hom}\bigl(\mathcal C^o,(\operatorname{Ens})\bigr)
\]
from $\mathcal C$ to the category of covariant functors from the dual
$\mathcal C^o$ of $\mathcal C$ to the category of sets. Recall:

\result{PROPOSITION}{1.1} The functor $h$ is \emph{fully faithful}; in other
words, for every pair $X,X'$ of objects of $\mathcal C$, the natural map
\[
  \operatorname{Hom}(X,X')\longrightarrow
  \operatorname{Hom}(h_X,h_{X'})
\]
is bijective.

In particular, if a functor
$F\in\operatorname{Hom}(\mathcal C^o,(\operatorname{Ens}))$ is isomorphic to
a functor of the form $h_X$, then $X$ is determined up to a unique isomorphism.
One then says that $F$ is \emph{representable}. The preceding proposition says
that the canonical functor $h$ defines an equivalence of $\mathcal C$ with the
full subcategory of $\operatorname{Hom}(\mathcal C^o,(\operatorname{Ens}))$
formed by the representable functors. This fact underlies the notion of a
“solution of a Universal Problem”: such a problem always consists in asking
whether a given functor (contravariant as here, or covariant in the dual case)
from $\mathcal C$ to $(\operatorname{Ens})$ is representable.

\src{370}
Moreover, by the very definition of a product in a category [1], the functor
$h:X\longmapsto h_X$ commutes with products whenever they exist (and, more
generally, with finite or infinite projective limits, in particular with fiber
products, with the formation of “kernels” [2], etc., whenever they exist):
there is an isomorphism of functors
\[
  h_{X\times X'}\xrightarrow{\sim}h_X\times h_{X'}
\]
whenever $X\times X'$ exists; that is, there are bijections, functorial in $Y$,
\[
  h_{X\times X'}(Y)\xrightarrow{\sim}h_X(Y)\times h_{X'}(Y).
\]
In particular, giving a morphism
\[
  X\times X'\longrightarrow X''
\]
in $\mathcal C$ (that is, a “law of composition” in $\mathcal C$ among
$X,X',X''$) is equivalent to giving a morphism
\[
  h_{X\times X'}=h_X\times h_{X'}\longrightarrow h_{X''};
\]
that is, to giving for every $Y\in\mathcal C$ a law of composition of sets
\[
  h_X(Y)\times h_{X'}(Y)\longrightarrow h_{X''}(Y)
\]
such that, for every morphism $Y\longrightarrow Y'$ in $\mathcal C$, the
system of set maps
\[
  h_{X^{(i)}}(Y)\longleftarrow h_{X^{(i)}}(Y')
  \qquad(\text{for }i=0,1,2)
\]
is a morphism for the two laws of composition belonging to $Y$ and $Y'$. In
this way one sees that a structure of “$\mathcal C$-group,”
“$\mathcal C$-ring,” etc. on an object $X$ of $\mathcal C$ is expressed most
conveniently (in theory as well as in practice) by saying that, for every
$Y\in\mathcal C$, the set $h_X(Y)$ carries a group law (respectively a ring
law, etc.) in the usual sense, and that the maps
$h_X(Y)\longleftarrow h_X(Y')$ corresponding to morphisms
$Y\longrightarrow Y'$ are group homomorphisms (respectively ring
homomorphisms, etc.). This is, for example, the most intuitive and convenient
way to define the various classical groups $G_a$, $G_m$,
$\operatorname{Gl}(n)$, etc. over an arbitrary base prescheme $S$ and to write
the classical relations among these groups; to put a “vector bundle” structure
on the affine scheme $V(\mathcal F)$ over $S$ defined by a quasi-coherent
sheaf $\mathcal F$; and to define and study the associated flag varieties
(Grassmannians, projective bundles), etc. The general yoga is simply to
identify, by means of the canonical functor $h$, the objects of $\mathcal C$
with particular contravariant functors, namely the representable functors,

\src{371}
from $\mathcal C$ to the category of sets.

The usual procedure of reversing arrows, which is required, for example, in
the case of affine schemes in passing from geometric language to the language
of commutative algebra, leads one to dualize the preceding considerations and,
in particular, also to introduce representable covariant functors
$\mathcal C\longrightarrow(\operatorname{Ens})$, that is, functors of the
form
\[
  Y\longmapsto\operatorname{Hom}(X,Y)=h'_X(Y).
\]

\subsection*{2. Pro-representable functors and pro-objects}
\addcontentsline{toc}{subsection}{2. Pro-representable functors and pro-objects}

Let $\widetilde X=(X_i)_{i\in I}$ be a projective system of objects of
$\mathcal C$. It has an associated covariant functor
\[
  h'_{\widetilde X}=\varinjlim_i h'_{X_i},
\]
or, more explicitly,
\[
  h'_{\widetilde X}(Y)=\varinjlim_i h'_{X_i}(Y)
  =\varinjlim_i\operatorname{Hom}(X_i,Y),
\]
from $\mathcal C$ to $(\operatorname{Ens})$. A functor from $\mathcal C$ to
$(\operatorname{Ens})$ that is isomorphic to a functor of this type with $I$
filtered is called \emph{pro-representable}. By the preceding subsection,
these are exactly the functors isomorphic to filtered inductive limits of
representable functors. Let
$\widetilde X'=(X'_j)_{j\in J}$ be a second filtered projective system in
$\mathcal C$ (constructed on a second filtered preordered index set $J$). One
readily verifies the canonical bijection (generalizing Proposition 1.1)
\[
  \operatorname{Hom}(h'_{\widetilde X'},h'_{\widetilde X})
  =\varprojlim_j\varinjlim_i\operatorname{Hom}(X_i,X'_j).
\]
This leads one to introduce the category $\operatorname{Pro}(\mathcal C)$ of
pro-objects of $\mathcal C$. Its objects are the projective systems of objects
of $\mathcal C$ (on variable filtered preordered index sets), and if
$\widetilde X=(X_i)_{i\in I}$ and
$\widetilde X'=(X'_j)_{j\in J}$ are two such objects, one puts
\[
  \operatorname{ProHom}(\widetilde X,\widetilde X')
  =\varprojlim_j\varinjlim_i\operatorname{Hom}(X_i,X'_j),
\]
the composition of pro-homomorphisms being evident. By construction,
$\widetilde X\longmapsto h'_{\widetilde X}$ may be regarded as a contravariant
functor of $\widetilde X$ and establishes an equivalence of the dual of the
category $\operatorname{Pro}(\mathcal C)$ of pro-objects of $\mathcal C$ with
the category of pro-representable covariant functors from $\mathcal C$ to
$(\operatorname{Ens})$. Of course, an object $X$ of $\mathcal C$ canonically
defines a pro-object, again denoted by $X$,

\src{372}
so that $\mathcal C$ is equivalent to a full subcategory of
$\operatorname{Pro}(\mathcal C)$. If
$\widetilde X=(X_i)_{i\in I}$ is any pro-object of $\mathcal C$, then (with the
preceding identification)
\[
  \widetilde X=\varprojlim_i X_i,
\]
where the projective limit is taken in $\operatorname{Pro}(\mathcal C)$
(because $h'_{\widetilde X}=\varinjlim_i h'_{X_i}$). Notice that even when the
projective limit of the $X_i$ exists in $\mathcal C$, it will generally not be
isomorphic to the projective limit $\widetilde X$ in
$\operatorname{Pro}(\mathcal C)$, as is already evident when $\mathcal C$ is
the category of sets. By definition,
$\varprojlim_{\mathcal C}X_i=L$ is defined by the requirement that the functor
of $Y\in\mathcal C$
\[
  \varprojlim_i\operatorname{Hom}_{\mathcal C}(Y,X_i)
  =\operatorname{Hom}_{\operatorname{Pro}(\mathcal C)}(Y,\widetilde X),
\]
with values in $(\operatorname{Ens})$, be representable by $L$, that is, be
isomorphic to $\operatorname{Hom}_{\mathcal C}(Y,L)$. Consequently,
$\varprojlim_{\mathcal C}X_i$ is already defined in terms of the pro-object
$\widetilde X$, and depends functorially on $\widetilde X$ whenever it is
defined; there is therefore no inconvenience in denoting it by
$\varprojlim_{\mathcal C}(\widetilde X)$. If projective limits always exist in
$\mathcal C$, then $\varprojlim_{\mathcal C}(\widetilde X)$ is a functor from
$\operatorname{Pro}(\mathcal C)$ to $\mathcal C$, and there is a canonical
homomorphism of functors
$\varprojlim_{\mathcal C}(\widetilde X)\longrightarrow\widetilde X$.

Every functor, say covariant to fix ideas,
\[
  F:\mathcal C\longrightarrow\mathcal C',
\]
extends in the evident way to a functor
\[
  \operatorname{Pro}(F):\operatorname{Pro}(\mathcal C)
  \longrightarrow\operatorname{Pro}(\mathcal C').
\]
It follows that, if projective limits always exist in $\mathcal C'$, then $F$
also canonically defines the composite functor
\[
  \overline F:\operatorname{Pro}(\mathcal C)\longrightarrow\mathcal C',
  \qquad
  \overline F=\varprojlim_{\mathcal C'}\circ\operatorname{Pro}(F),
\]
which sends $\widetilde X=(X_i)_{i\in I}$ to
$\varprojlim_{\mathcal C'}F(X_i)$.

A pro-object $\widetilde X$ is called a \emph{strict pro-object} if it is
isomorphic to a pro-object $(X_i)_{i\in I}$ whose transition morphisms
$X_i\longrightarrow X_j$ are epimorphisms; a functor defined by such an object
is called \emph{strictly pro-representable}. One may then further require $I$
to be filtered and ordered and every epimorphism $X_i\longrightarrow X'$ to be
equivalent to an epimorphism

\src{373}
$X_i\longrightarrow X_j$ for a suitable $j\in I$ (uniquely determined by this
condition). Under these conditions the projective system $(X_i)_{i\in I}$ is
determined up to a unique isomorphism (in the usual sense of isomorphisms of
projective systems). It follows that in $\operatorname{Pro}(\mathcal C)$ the
projective limit of a projective system of strict pro-objects
$\widetilde X^{(\alpha)}$ always exists and, with the notation above for $F$
and $\overline F$, one has
\[
  \overline F\left(\varprojlim_\alpha\widetilde X^{(\alpha)}\right)
  =\varprojlim_{\alpha,\mathcal C'}
  \overline F\left(\widetilde X^{(\alpha)}\right).
\]
In particular, if every pro-object of $\mathcal C$ is strict (cf. the next
subsection), the extended functor $\overline F$ commutes with projective
limits.

\subsection*{3. Characterization of pro-representable functors}
\addcontentsline{toc}{subsection}{3. Characterization of pro-representable functors}

Let $\mathcal C$ and $\mathcal C'$ be two categories in which finite
projective limits (that is, limits relative to finite preordered sets, not
necessarily filtered) exist, or, equivalently, in which finite products and
finite fiber products exist. This implies in particular the existence of a
“right unit object” $e$ such that $\operatorname{Hom}(X,e)$ consists of one
element for every $X$. Let $F$ be a covariant functor from $\mathcal C$ to
$\mathcal C'$. The following conditions are equivalent:

\begin{enumerate}
\item[(i)] $F$ commutes with finite projective limits;
\item[(ii)] $F$ commutes with finite products and finite fiber products;
\item[(iii)] $F$ commutes with finite products, and for every exact diagram
\[
  X\longrightarrow X'\rightrightarrows X''
\]
in $\mathcal C$ ([3], A, Definition 2.1), the diagram obtained by applying $F$,
\[
  F(X)\longrightarrow F(X')\rightrightarrows F(X''),
\]
is exact.
\end{enumerate}

One then says that $F$ is \emph{left exact}.

Henceforth we assume that finite projective limits exist in $\mathcal C$. It
is then immediate from the definitions that a representable functor is left
exact, and, on passing to the limit, that a pro-representable functor is left
exact.

For a converse, let
\[
  F:\mathcal C\longrightarrow(\operatorname{Ens})
\]
be a covariant functor, and let $X\in\mathcal C$ and $\xi\in F(X)$. We say
that $\xi$ (or the pair \src{374}$(X,\xi)$) is \emph{minimal} if, for every
$X'\in\mathcal C$, every $\xi'\in F(X')$, and every strict monomorphism ([3],
A, 2) $u:X'\longrightarrow X$ such that $\xi=F(u)(\xi')$, the morphism $u$ is
an isomorphism. On the other hand, a pair $(X,\xi)$ is said to
\emph{dominate} $(X'',\xi'')$ ($\xi\in F(X)$, $\xi''\in F(X'')$) if there is
a morphism $v:X\longrightarrow X''$ such that $\xi''=F(v)(\xi)$. If $\xi$ is
minimal and $F$ is left exact, this morphism $v$ is unique; if $\xi''$ is
minimal, $v$ is surjective. The following proposition follows readily:

\result{PROPOSITION}{3.1} The functor $F$ is strictly pro-representable if and
only if it satisfies the following two conditions:
\begin{enumerate}
\item[(i)] $F$ is left exact;
\item[(ii)] every pair $(X,\xi)$, with $\xi\in F(X)$, is dominated by a
minimal pair.
\end{enumerate}

The latter condition is vacuous if every object of $\mathcal C$ is Artinian
(take a subobject $X'$ of $X$ minimal among those for which there exists
$\xi'\in F(X')$ whose image is $\xi$). Hence:

\result{COROLLARY}{} Let $\mathcal C$ be a category whose objects are Artinian
and in which finite projective limits exist. Then the pro-representable
functors from $\mathcal C$ to $(\operatorname{Ens})$ are exactly the left exact
functors, and they are in fact strictly pro-representable.

The latter fact also means that \emph{every pro-object of $\mathcal C$ is
strict}.

\subsection*{4. Example: groups of Galois type and pro-algebraic groups}
\addcontentsline{toc}{subsection}{4. Groups of Galois type and pro-algebraic groups}

If $\mathcal C$ is the category of ordinary finite groups, then
$\operatorname{Pro}(\mathcal C)$ is equivalent to the category of compact
totally disconnected topological groups. Groups of this type, and their
generalizations obtained by replacing ordinary finite groups with finite group
schemes over a given base prescheme (for example, finite algebraic groups over
a field $k$), will serve as absolute and relative fundamental, homotopy, and
homology groups for preschemes. In all these examples, the corollary to
Proposition 3.1 applies, and $\pi_1$ must indeed be defined by means of the
associated functor [2]. The same is true when one starts from the category of
algebraic or quasi-algebraic groups over a field (or, more generally, over a
Noetherian prescheme): one obtains Serre's “pro-algebraic groups” [4].

\subsection*{5. Example: “formal varieties”}
\addcontentsline{toc}{subsection}{5. Example: formal varieties}

Let $\Lambda$ be a Noetherian ring, and let $\mathcal C$ be the category of
$\Lambda$-algebras $A$ that are Artinian finite-type modules over $\Lambda$
(or, more briefly,

\src{375}
Artinian $\Lambda$-algebras). The hypotheses of the corollary to Proposition
3.1 are satisfied. Here the category $\operatorname{Pro}(\mathcal C)$ is
equivalent to the category of topological algebras $O$ over $\Lambda$ that are
isomorphic to topological projective limits
\[
  O=\varprojlim_i O_i
\]
of algebras $O_i\in\mathcal C$; that is, their topology is linear, separated,
and complete, and for every open ideal $J_i$ of $O$, the quotient $O/J_i$ is
an Artinian algebra over $\Lambda$. The functor
$\mathcal C\longrightarrow(\operatorname{Ens})$ associated with such an
algebra is precisely
\begin{align*}
  F(A)=h'_O(A)
  &=\{\text{continuous homomorphisms of topological $\Lambda$-algebras }
       O\longrightarrow A\}\\
  &=\varinjlim_i
    \operatorname{Hom}_{\Lambda\text{-algebras}}(O_i,A).
\end{align*}

Notice also that $\mathcal C$ is essentially the product of the analogous
categories corresponding to the completed local rings of the
$\Lambda_{\mathfrak m}$ for the maximal ideals $\mathfrak m$ of $\Lambda$; if
desired, one may therefore restrict to the case in which $\Lambda$ is such a
complete local ring. In every case, $O$ decomposes canonically as the
topological product of local components corresponding to the “points” of the
formal scheme [2] defined by $O$. Such a point is defined by an object $\xi$
of some $F(K)$, where $K\in\mathcal C$ is a field (for example, the residue
field of the local component in question); two pairs $(\xi,K)$ and
$(\xi',K')$ define the same point if and only if they are dominated by a common
$(\xi'',K'')$, or, equivalently, if they dominate a common
$(\xi''',K''')$. (If the $\Lambda/\mathfrak m$ are algebraically closed, it
suffices to take the disjoint union of the $F(\Lambda/\mathfrak m)$.)

It is important to give conditions under which the local component $O_\xi$ of
$O$ corresponding to $\xi\in F(K)$ is a Noetherian ring. When $\Lambda$ is a
complete local ring (Noetherian, as we recall), this is equivalent to saying
that $O_\xi$ is isomorphic to a quotient of a formal power series ring
$\Lambda[[t_1,\ldots,t_n]]$ over $\Lambda$. To give such a criterion, for
every ring $A$ introduce the $A$-algebra $I_A$ of the “dual numbers” of $A$ by
\[
  I_A=A[t]/t^2A[t].
\]
Let $\varepsilon:I_A\longrightarrow A$ be the augmentation homomorphism. If
$A\in\mathcal C$, it defines a map
\[
  F(\varepsilon):F(I_A)\longrightarrow F(A),
\]

\src{376}
and, using the fact that $F$ is left exact, one intrinsically defines an
$A$-module structure on the subset
\[
  F(I_A,\xi)=F(\varepsilon)^{-1}(\xi)
\]
of $F(I_A)$ formed by the $\xi'\in F(I_A)$ that “reduce to $\xi$.” Using the
explicit form of $F$ in terms of $O$, one finds that this $A$-module is
identified with
\[
  \operatorname{Hom}_{\Lambda}(\mathfrak m_\xi/\mathfrak m_\xi^2,A),
\]
where $\mathfrak m_\xi$ is the kernel of the homomorphism
$\xi:O\longrightarrow A$; if $A$ is a field, this is the maximal ideal of the
local component $O_\xi$ of $O$. The following proposition is immediate:

\result{PROPOSITION}{5.1} Let $\xi\in F(K)$, where $K\in\mathcal C$ is a
field. The corresponding local component $O_\xi$ of $O$ is Noetherian if and
only if the set $F(I_K,\xi)$ of elements of $F(I_K)$ reducing to $\xi$ is a
finite-dimensional vector space over $K$. Under these conditions there is a
canonical isomorphism
\[
  F(I_K,\xi)=
  \operatorname{Hom}\!\left(
    \frac{\mathfrak m_\xi/\mathfrak m_\xi^2}
         {\operatorname{im}(\mathfrak n_\xi/\mathfrak n_\xi^2)},K
  \right),
\]
where $\mathfrak n_\xi$ is the maximal ideal of $\Lambda$, the kernel of
$\Lambda\longrightarrow K$.
\footnote{The printed formula is correct only when $\Lambda$ is a field. The
1962 erratum directs that in general
$\mathfrak m_\xi/\mathfrak m_\xi^2$ be replaced by its quotient by the image
of $\mathfrak n_\xi/\mathfrak n_\xi^2$.}
In particular, the dimension of the $K$-vector space $F(I_K,\xi)$ equals the
dimension of this corrected vector space over the field
$O_\xi/\mathfrak m_\xi=K(\xi)$.

Suppose that $O_\xi$ is Noetherian, and, to simplify the notation, suppose
that $\Lambda$ is complete local and $O=O_\xi$. We say that $O$ is
\emph{smooth over $\Lambda$} if $O$ is finite and étale over the completion of
the localization of $\Lambda[t_1,\ldots,t_n]$ at a maximal ideal inducing the
maximal ideal of $\Lambda$. When the residue extension of $O$ over $\Lambda$
is trivial (for example, when the residue field of $\Lambda$ is algebraically
closed), this is equivalent to saying that $O$ itself is isomorphic to such a
formal power series algebra.
\footnote{The 1962 erratum specifies that this definition is correct only when
the residue extension $k'/k$ is separable; for the general case it refers to
SGA III, 1.1.}
Finally, if $O$ is no longer assumed Noetherian, we still say that $O$ is
smooth over $\Lambda$ if $O$ is isomorphic to a topological projective limit
of quotient $\Lambda$-algebras that are Noetherian and smooth over $\Lambda$
in the preceding sense. This generalizes at once to the case in which
$\Lambda$ and $O$ are no longer assumed local. We then have:

\result{PROPOSITION}{5.2} The algebra $O$ is smooth over $\Lambda$ if and only
if the associated functor $F$ sends epimorphisms to epimorphisms.

Thus, for every surjective homomorphism $A\longrightarrow A'$ in
$\mathcal C$, the map $F(A)\longrightarrow F(A')$ is surjective. It is of
course enough to verify this condition when $A$ is local and, proceeding step
by step, when the ideal of $A$ that is the kernel of $A\longrightarrow A'$ is
annihilated by the maximal ideal of $A$. In practice, this reduces matters to
checking that a certain obstruction, related to infinitesimal invariants

\src{377}
of the situation giving rise to the functor $F$, vanishes, a problem of a
cohomological nature.

To conclude, let us say a few words in the preceding context about
\emph{rings of definition}. Let $F$ still be a functor from $\mathcal C$ to
$(\operatorname{Ens})$, pro-represented by a topological $\Lambda$-algebra
$O$. Then for every $A\in\mathcal C$ and every $\xi\in F(A)$, there is a
smallest subring $A'$ of $A$ such that $\xi$ is the image of an element $\xi'$
of $F(A')$ (which is then uniquely determined). Indeed, it suffices to
interpret $\xi$ by means of a homomorphism from $O$ to $A$ and to take for
$A'$ the image of $O$ under this homomorphism. One then says that $A'$ is the
ring of definition of the object $\xi\in F(A)$. If $u:A\longrightarrow B$ is
an algebra homomorphism and $\eta=F(u)(\xi)$, then the ring of definition of
$\eta$ is the image under $u$ of the ring of definition of $\xi$. When one
starts from a functor $F$ from $\mathcal C$ to $(\operatorname{Ens})$, the
existence of rings of definition and the properties just given are, up to
minor qualifications, equivalent to $F$ being pro-representable; thus they are
usually far from trivial.

\fsection{B}{The two existence theorems}

Retain the notation of A, Section 5, and start with a covariant functor
\[
  F:\mathcal C\longrightarrow(\operatorname{Ens}).
\]
We seek manageable criteria for $F$ to be pro-representable, that is,
expressible by means of a $\Lambda$-algebra $O$ as above. By the corollary to
A, Proposition 3.1, this is equivalent to $F$ being left exact. In the present
state of the descent technique (cf. the questions posed in [3], p. 9), this
criterion cannot be directly verified in this form in the most important
cases, and criteria that appear less demanding are needed.

\result{THEOREM}{1} The functor $F$ is pro-representable if and only if it
satisfies the following two conditions:
\begin{enumerate}
\item[(i)] $F$ commutes with finite products;
\item[(ii)] for every algebra $A\in\mathcal C$ and every homomorphism
$A\longrightarrow A'$ in $\mathcal C$ such that the diagram
\[
  A\longrightarrow A'\rightrightarrows A'\otimes_A A'
\]
is exact ([3], A, Definition 1.2), the transformed diagram
\[
  F(A)\longrightarrow F(A')\rightrightarrows F(A'\otimes_A A')
\]

\src{378}
is exact.
\end{enumerate}
Moreover, it is enough to verify (ii) when $A$ is local and, in addition, one
of the following two cases holds:
\begin{enumerate}
\item[(a)] $A'$ is a free module over $A$;
\item[(b)] the quotient module $A'/A$ is an $A$-module of length 1.
\end{enumerate}

The proof of this theorem is rather delicate and cannot be sketched here. We
merely note that it rests essentially on a study of equivalence relations (in
the categorical sense) in the spectrum of an Artinian algebra, a study that
still raises several problems whose solution seems indispensable for the
further development of the theory.

In applications, condition (i) is always trivial to verify. The verification
of (ii) splits in two: the case in which $A'$ is a free $A$-module belongs to
the descent technique along flat morphisms ([3], Theorems 1, 2, and 3) and
presents no difficulty. To treat case (b), we use the following result:

\result{THEOREM}{2} Let $A$ be an Artinian local ring with maximal ideal
$\mathfrak m$, and let $A'$ be an $A$-algebra containing $A$ such that
$\mathfrak mA'\subset A$ and
\[
  A\longrightarrow A'\rightrightarrows A'\otimes_A A'
\]
is exact (as is the case, in particular, if $A'/A$ is an $A$-module of length
1). Let $\mathcal F$ be the fibered category ([3], A, Definition 1.1) of flat
quasi-coherent sheaves on variable preschemes. Then the morphism
$\operatorname{Spec}(A')\longrightarrow\operatorname{Spec}(A)$ is a strict
$\mathcal F$-descent morphism ([3], A, Definition 1.7).

In other words, giving a flat $A$-module $M$ is completely equivalent to
giving a flat $A'$-module $M'$ together with an
$A'\otimes_AA'$-isomorphism from $M'\otimes_AA'$ onto $A'\otimes_AM'$ that
satisfies the usual transitivity condition for a descent datum (loco citato).

Theorem 2 is proved by first showing that
\[
  H^i(A'/A,G_a)=0\qquad\text{for }i\geq 1
\]
([3], A, 4(e)); the hypothesis $\mathfrak mA'\subset A$ allows one to reduce
readily to the case in which $A$ is a field $k$ (namely $A/\mathfrak m$). One
then applies the equivalences noted in [3], page 16.

\fsection{C}{Applications to some special cases}

\subsection*{1. General remarks on functors represented by preschemes}
\addcontentsline{toc}{subsection}{1. General remarks on functors represented by preschemes}

Let $S$ be a locally Noetherian prescheme. A prescheme $X$ over $S$ is called
\src{379}locally of finite type over $S$ if, for every $x\in X$ mapping to
$y\in S$,\footnote{The printed text has “$y\in Y$” here and uses $Y$ again at
the end of the paragraph. Since the morphism under consideration is
$X\longrightarrow S$, both occurrences are read as $S$.} there is an affine
neighborhood of $y$, with ring $A$, and an affine neighborhood of $x$ over it,
with ring $B$, such that $B$ is an $A$-algebra of finite type. One finds many
important examples of preschemes locally of finite type over $S$ that are not
of finite type over $S$, as solutions of classical universal problems. Thus it
is important to be able to regard the Picard scheme of a curve as the union of
infinitely many connected components (which must not be confused with the
connected component of the identity element, i.e. the “Picard variety”). It is
then sometimes convenient to work in the category $\mathcal C$ of preschemes
locally of finite type over $S$, and there examine the representability of a
contravariant functor $F$. The principal aim of these exposés is to develop a
general technique for recognizing whether such a functor $F$ is representable,
and for studying the properties of the corresponding $S$-prescheme $X$ by
means of those of $F$. Let us mention in passing that this study produces
nonpathological examples of preschemes over $S$ that are not separated over
$S$, notably as “Picard preschemes” of excellent $S$-schemes. Preschemes that
are not schemes must therefore not be banished from algebraic geometry.

Let $X$ be a prescheme locally of finite type over $S$, and let
\[
  F(Y)=\operatorname{Hom}_S(Y,X)
\]
be the associated contravariant functor. We may consider the restriction $F_0$
of $F$ to the subcategory $\mathcal C_0$ of $\mathcal C$ formed by the
preschemes $Y$ over $S$ that are Artinian and finite over $S$: when
$S=\operatorname{Spec}(\Lambda)$, $\mathcal C_0$ is therefore the dual of the
category of Artinian $\Lambda$-algebras considered in B. If
$Y=\operatorname{Spec}(A)$, where $A$ is an Artinian local ring, then $Y$
consists of a single point $y$ lying over a closed point $s$ of $S$, and an
$S$-homomorphism from $Y$ to $X$ (i.e. an element of $F(Y)$) is specified by a
point $x\in X$ over $s$ and an $\mathcal O_{s,S}$-homomorphism from
$\mathcal O_{x,X}$ to $A$. If such a homomorphism exists, then $x$ is
necessarily a closed point of $X$ (its residue field being algebraic over that
of $s$). This shows that the restriction $F_0$ of $F$ to the “Artinian
$\Lambda$-algebras” is pro-representable, and is represented by the topological
$\Lambda$-algebra whose local components are the completions
$\widehat{\mathcal O}_{x,X}$ of the local rings of $X$ at the closed points
$x$ of $X$ lying over closed points of $S$. Thus knowledge of $F_0$ alone gives
precise information about the structure of $X$, namely the structure of the
completions of its local rings at the indicated points. Note that even when
$S$

\src{380}
is the spectrum of an algebraically closed field, it is only through the
systematic consideration of “varieties” $Y$ for which $\mathcal O_Y$ may have
nilpotent elements, and in particular by working with spectra of Artinian local
rings, that one arrives at the “correct formulation” of classical universal
problems and grasps their “infinitesimal” aspect.

When one starts from a functor $F$ given in advance and wants to decide whether
it is representable, the study of $F_0$ (using Theorems 1 and 2) gives nearly
complete indications. It may happen, as it frequently does (for example, by
testing simply the nature of the sets $F(I_K,\xi)$ and their functorial
behavior; cf. A), that $F_0$ is already found not to be pro-representable. This
explains the failure of the attempts made thus far to define, in a more or less
natural way, moduli varieties classifying vector bundles of rank $>1$. Or one
may verify that $F_0$ is representable but that the vector spaces
$F(I_K,\xi)$ are not finite-dimensional, in which case one must be content
with the “formal” solution. Finally, $F_0$ may be represented by a product of
complete Noetherian local rings. This gives very strong grounds for believing
that $F$ itself is representable and, together with analogous but more global
properties that we intend to develop later, will doubtless suffice to imply
that it actually is. Lastly, one encounters interesting geometric problems
(see paragraphs 4 and 5 below) for which only the functor $F_0$ is available
(and does not arise from any “global” functor $F$), and one will count oneself
fortunate if a “formal moduli scheme” can be associated with it.

To conclude these general remarks, let us indicate how scheme theory explains
apparent anomalies such as the Igusa surface $V$, whose “Picard variety” $P$
is a point even though
\[
  H^1(V,\mathcal O_V)\ne 0.
\]
In this case $P$ is a “purely infinitesimal” group not reduced to the identity
group; that is, it is defined by a local algebra $O$ of finite rank over the
base field $k$, equipped with a diagonal map corresponding to the
multiplicative structure of $P$. If $\mathfrak m$ is the maximal ideal of $O$,
the dual of $\mathfrak m/\mathfrak m^2$ is canonically isomorphic to
$H^1(V,\mathcal O_V)$ (cf. paragraph 3 below). Only when the Picard group is an
algebraic group in the classical sense, i.e. smooth over the base field $k$,
is the dimension of $H^1(V,\mathcal O_V)$ (always equal to that of
$\mathfrak m/\mathfrak m^2$) equal to the dimension of the Picard group.

\subsection*{2. Schemes $\underline{\operatorname{Hom}}_S(X,Y)$,
$\Pi_{X/S}Z$, $\underline{\operatorname{Aut}}_S(X)$, etc.}
\addcontentsline{toc}{subsection}{2. Hom, Pi, and Aut schemes}

Let $X,Y$ be two preschemes over $S$; for every prescheme $T$ over $S$,

\src{381}
put $X_T=X\times_ST$, $Y_T=Y\times_ST$, and consider the set
\[
  F(T)=\operatorname{Hom}_T(X_T,Y_T)
      =\operatorname{Hom}_S(X_T,Y)
      =\operatorname{Hom}_S(X\times_ST,Y)
\]
as a contravariant functor of $T$. If it is representable, the prescheme over
$S$ that represents it is denoted by
$\underline{\operatorname{Hom}}_S(X,Y)$; thus there is a functorial
isomorphism
\[
  \operatorname{Hom}_S
  \bigl(T,\underline{\operatorname{Hom}}_S(X,Y)\bigr)
  \xrightarrow{\sim}
  \operatorname{Hom}_S(T\times_SX,Y).
\]
There are many variants of this universal problem whose solution reduces to
it: the prescheme of $S$-automorphisms of an $S$-prescheme $X$ (which will be a
group prescheme), the prescheme of $S$-homomorphisms from one $S$-group
prescheme to another (which will be a commutative group prescheme if the second
group scheme is commutative), etc. One may also generalize the definition of
$\underline{\operatorname{Hom}}_S(X,Y)$ by considering a prescheme $Z$ over
the prescheme $X$ over $S$, and the functor
\[
  F(T)=\operatorname{Hom}_{X_T}(X_T,Z_T)
\]
(the set of “sections” of $Z_T$ over $X_T$). If this functor is representable,
the $S$-prescheme representing it will be denoted $\Pi_{X/S}Z$; hence, by
definition, there is a functorial isomorphism
\[
  \operatorname{Hom}_S(T,\Pi_{X/S}Z)
  =\operatorname{Hom}_{X_T}(X_T,Z_T).
\]
Taking $Z=Y\times_SX$ recovers $\underline{\operatorname{Hom}}_S(X,Y)$. These
definitions give a formulary for the newly introduced preschemes that is as
trivial as it is useful; we shall not give it here, since it holds in every
category in which products and fiber products exist. More serious is the
question of the existence of schemes of the form
$\underline{\operatorname{Hom}}_S(X,Y)$. One first observes that, for fixed
$X$, $\underline{\operatorname{Hom}}_S(X,Y)$ can exist for every $Y$ over $S$
(or $\Pi_{X/S}Z$ for every $Z$ over $X$) only if $X$ is flat over $S$.
Moreover, one is convinced that it is reasonable to expect a solution for
sufficiently general $Y$ only if $X$ is also proper over $S$. It appears,
furthermore, that these conditions suffice for the existence of
$\underline{\operatorname{Hom}}_S(X,Y)$ and $\Pi_{X/S}Z$, provided that, when
necessary, a “quasi-projective” hypothesis is imposed on $Y/S$ or on $Z/X$,
respectively. This can in any case be verified in a fair number of situations
(for example, when $Y$ is affine over $S$, or, by direct elementary
constructions, when $X$ is finite over $S$). Theorems 1 and 2 give the
following:

\src{382}
\result{PROPOSITION}{2.1} Let $\Lambda$ be a Noetherian ring and let $X,Y$ be
arbitrary preschemes over $\Lambda$. Consider the functor
\[
  F(A)=\operatorname{Hom}_A(X_A,Y_A)
\]
on the category $\mathcal C_0$ of Artinian $\Lambda$-algebras. If $X$ is flat
over $\Lambda$, this functor is pro-representable.

Moreover, for every $A\in\mathcal C_0$ and $\xi\in F(A)$, there is a canonical
isomorphism
\[
  F(I_A,\xi)=H^0\!\left(
    X_A,
    \underline{\operatorname{Hom}}_{\mathcal O_{X_A}}
    \bigl(\xi^*(\Omega^1_{Y_A/A}),\mathcal O_{X_A}\bigr)
  \right),
\]
where $\Omega^1_{Y_A/A}$ is the sheaf of Kähler $1$-differentials of $Y_A$
over $A$. Taking $A$ to be a field and using A.5.1 and the finiteness theorem
of [2], one obtains the following corollary:

\result{COROLLARY}{} Suppose that $X$ is flat and proper over $S$ and that $Y$
is locally of finite type over $S$. Then $F$ is pro-representable, and the local
components of the corresponding topological $\Lambda$-algebra are Noetherian
rings.

\result{REMARKS}{} The problems considered in this number, and many others,
were commonly studied in classical algebraic geometry by means of the “Chow
coordinates” of cycles in projective space, which allowed these cycles to be
regarded as points of suitable projective varieties. This procedure, and the
use of Chow coordinates in general, seems irredeemably inadequate from the
scheme-theoretic viewpoint because it destroys nilpotent elements in parameter
varieties and, in particular, is not suited to a satisfactory study of
infinitesimal variations of cycles (quite apart from its nonintrinsic nature,
which is tied to projective space). Unfortunately, Chow coordinates were the
only language used by many algebraic geometers to study families of varieties
or cycles. This seems to have seriously obstructed the clarification of these
notions, despite the method's undoubted (and probably temporary) technical
value. To obtain the analogue of Chow varieties in scheme theory, one is led
to the following universal problem. Let $X$ be a prescheme over $S$. For every
prescheme $T$ over $S$, let $F(T)$ be the set of closed subpreschemes of
$X_T=X\times_ST$ that are flat over $T$. One seeks to represent this functor
of $T$ by a prescheme over $S$. More generally, one may take a quasi-coherent
sheaf $G$ on $X$ and let $F(T)$ be the set of quotient sheaves

\src{383}
of $G_T$ that are flat over $T$. It appears that the problem has a solution,
with a scheme $C$ locally of finite type over $S$, when $X$ is proper over the
locally Noetherian prescheme $S$ and $G$ is coherent. In any case, assuming
only that $S$ is locally Noetherian, the restriction of $F$ to the “Artinian
$S$-algebras” is pro-representable; and if $X$ is also proper over $S$ and $G$
is coherent, the local components of the corresponding topological ring $O$
are Noetherian. Of course, even after the existence of the “Chow scheme” of
$X$ over $S$ has been proved, one will still have to find a decomposition into
disjoint open subsets $C_i$ (corresponding to fixing continuous invariants such
as the degree and dimension of the varying cycles) that are of finite type,
not merely locally of finite type, over $S$; determine the relation between
this scheme and the classical Chow varieties; and determine when a $C_i$ is
even projective, or at least quasi-projective, over $S$.\footnote{The 1962
erratum records that the projective case is completely solved by Hilbert
schemes (TDTE IV), while examples of Nagata and Hironaka show that without a
projectivity hypothesis these functors need not be representable, even for
zero-dimensional subvarieties of a nonsingular complete threefold; this is
connected with the possible nonexistence of the symmetric square.}

\subsection*{3. Picard schemes}
\addcontentsline{toc}{subsection}{3. Picard schemes}

Let $f:X\longrightarrow S$ be an $S$-prescheme.\footnote{For a fuller study,
the 1962 erratum refers to TDTE V.} Consider the multiplicative
sheaf $\mathcal O_X^*$ of units in the structure sheaf of $X$, and the group
\[
  P(X/S)=H^0\bigl(S,R^1f_*(\mathcal O_X^*)\bigr),
\]
called the relative Picard group of $X/S$. An element of this group is thus
specified by an open covering $(U_i)$ of $S$ and an invertible sheaf $L_i$ on
each $f^{-1}(U_i)$, in such a way that for every pair $(i,j)$ the restriction
$L_i|f^{-1}(U_i\cap U_j)$ is isomorphic to
$L_j|f^{-1}(U_i\cap U_j)$, at least locally over $U_i\cap U_j$ (that is, the
two sheaves are “equivalent” in the sense of [3], B, 4). If $X/S$ has a
section, then $P(X/S)$ is simply the set of classes of invertible sheaves over
$X/S$, up to “equivalence” (loco citato). Now put, for every $T$ over $S$,
\[
  F(T)=P(X_T/T).
\]
This gives a contravariant functor of $T$, which may be called the Picard functor of
$X/S$. If this functor is representable, the prescheme over $S$ that represents
it is called the Picard prescheme of $X/S$ and is denoted $\underline P(X/S)$.
Thus there is an isomorphism of functors
\[
  \operatorname{Hom}_S\bigl(T,\underline P(X/S)\bigr)
  \xrightarrow{\sim}P(X_T/T).
\]
The formation of Picard preschemes is compatible with base extension. In
particular, the Picard preschemes of the fibers of $X$ over $S$ (which are
preschemes over the residue fields $K(s)$ of the points $s\in S$) are the
fibers of $\underline P(X/S)$. Of course, since $P(X_T/T)=F(T)$ is a
commutative group, Picard preschemes are group preschemes. Note also that the
generalized Jacobians of Rosenlicht are simply the connected components of the
identity in the Picard schemes of complete curves that may have singularities;
this should make most of their properties evident once existence has been
proved.

\result{REMARK}{} The definition adopted here is reasonable only when every
point of $S$\footnote{The printed text has “$Y$”; no base $Y$ is defined here,
and the condition concerns the base $S$.} has an open neighborhood $U$ over
which $X$ has a section. In the general case the definition of the Picard
functor must be modified slightly in order still to obtain an existence
theorem.

Here the plausible conditions for the existence of a Picard prescheme are the
following: $X$ is proper and flat over $S$, $f_*(\mathcal O_X)=\mathcal O_S$,
and $X$ locally has a section over $S$. The last condition enters naturally in
the application of the descent technique: it eliminates the automorphisms of
an invertible sheaf $L$ on $X$ by equipping it with a marked section over the
section $s$ ([3], B, 4). In particular, one obtains:

\src{384}
\result{PROPOSITION}{3.1} Suppose that $X$ is flat over
$S=\operatorname{Spec}(\Lambda)$, where $\Lambda$ is Noetherian, and that for
every $T$ of finite type over $S$ one has
\[
  f_{T*}(\mathcal O_{X_T})=\mathcal O_T
\]
(if $f$ is proper and separable,\footnote{The 1962 erratum explicitly prescribes
the hypothesis that $f$ be “proper and separable.”} or if $S$ is the spectrum
of a field, Künneth's theorem shows that the last condition is equivalent to
$f_*(\mathcal O_X)=\mathcal O_S$). Then the Picard functor of $X/S$ on the
category of Artinian $\Lambda$-algebras is pro-representable.

Moreover, here one has
\[
  F(I_A,\xi)=H^1(X_A,\mathcal O_{X_A});
\]
in particular:

\result{COROLLARY}{} If $X$ is proper over $S$, then the local components of
the topological $\Lambda$-algebra $O$ corresponding to the Picard functor are
Noetherian.

\result{REMARKS}{} The definitions and results of this number may be
generalized to the classification of principal bundles on $X$ whose structure
group is a group scheme $G$ over $S$ that is affine, flat over $S$, and
commutative. If $G$ is not commutative, then the adjoint group bundle of a
principal bundle (whose sections are the automorphisms of the principal
bundle) is no longer trivial, and Proposition 3.1 is no longer valid as stated.
The universal problem can nevertheless be modified so as still to obtain a
solution (at least, for the moment, in formal geometry). The golden rule to
retain, in the context of this \src{385}number and those that follow, and whenever one
seeks “moduli schemes” for classes of objects defined only up to isomorphism,
is always this: eliminate the possible automorphisms of the objects to be
classified by introducing auxiliary structures when necessary (marked points
or elements of sections, fixed differential forms, etc.), chosen innocuously
enough not to alter the original problem substantially.

\subsection*{4. Formal moduli of a variety}
\addcontentsline{toc}{subsection}{4. Formal moduli of a variety}

Let $\Lambda$ be a Noetherian local ring with residue field $k$ (most often
$\Lambda$ will equal $k$, or will be a Cohen $p$-ring), and let $X_0$ be a
prescheme over $k$. For every Artinian local $\Lambda$-algebra $A$, consider
the set $F(A)$ of isomorphism classes of $A$-preschemes $X$ flat over $A$,
equipped with an isomorphism
\[
  (*)\qquad
  X\otimes_A\kappa(A)\xleftarrow{\sim}
  X_0\otimes_k\kappa(A),
\]
where $\kappa(A)$ is the residue field of $A$; of course, isomorphisms between
such flat $A$-preschemes must respect the preceding isomorphism included in the
structure. If $A$ is an Artinian $\Lambda$-algebra not necessarily local, with
local components $A_i$, define $F(A)$ to be the product of the $F(A_i)$. Thus
$F$ becomes a multiplicative functor of $A$, which may be called the moduli
functor for $X_0$ (and $\Lambda$). If this functor is representable, it
corresponds to a local topological $\Lambda$-algebra $O$ with residue field
$k$, and the formal spectrum of $O$ is called the formal moduli scheme for
$X_0$ (and $\Lambda$); cf. [2] for some details on this notion.

Here, for applying the descent technique, the “finite” automorphisms of $X_0$
are harmless, since they have no bearing on the existence of automorphisms (in
the sense specified above) of the $A$-preschemes $X$. When $A$ is not a field,
the necessary and sufficient condition for $X$ to have no nontrivial
$A$-automorphism is
\[
  H^0(X_0,\Theta_{X_0/k})=0,
\]
where $\Theta_{X_0/k}$ denotes the sheaf of $k$-derivations (the tangent sheaf)
of $X_0$. One also finds readily (at least if $X_0$ is smooth over $k$) that
\[
  F(I_A,\xi)=H^1(X_A,\Theta_{X_A/A}).
\]
Hence one concludes, as usual:

\src{386}
\result{PROPOSITION}{4.1} Suppose that
$H^0(X_0,\Theta_{X_0/k})=0$. Then the formal moduli scheme for $X_0$ exists. If
$X_0$ is also proper over $k$, the formal moduli scheme is Noetherian.

\result{REMARKS}{} 1. When $X_0$ is not assumed smooth over $k$,
$F(I_A,\xi)$ is identified with an $A$-submodule of
\[
  \operatorname{Ext}^1_{\mathcal O_{P_A}}
  \bigl(P_A;\mathcal I_{X_A},\mathcal O_{X_A}\bigr),
\]
where $P_A=X_A\times_\Lambda X_A$, where $\mathcal O_{X_A}$ is regarded as a
coherent sheaf on $P_A$ by means of the diagonal morphism
$X_A\longrightarrow P_A$, and where $\mathcal I_{X_A}$ denotes the coherent
ideal sheaf on $P_A$ defined by the diagonal morphism. More precisely, a
straightforward globalization of Hochschild theory shows that the above
$\operatorname{Ext}^1$ is identified with the set of isomorphism classes of
flat sheaves of $I_A$-algebras $\mathcal O$ on $X_A$, equipped with an
augmentation isomorphism
\[
  \mathcal O\otimes_{I_A}A\longrightarrow\mathcal O_{X_A}
\]
(recall that $I_A=A[t]/(t^2)$). The submodule $F(I_A,\xi)$ is the one
corresponding to sheaves of commutative algebras. The smoothness hypotheses are
therefore not essential in moduli theory, as [2] had suggested.

2. Recall (loco citato) that, in particular, every algebraic curve $X_0$ that
is smooth and proper over $k$ has a formal moduli scheme, smooth over
$\Lambda$, of relative dimension $3g-3$ if its genus $g$ is at least $2$, and
$g$ if $g=0$ or $1$. The latter two cases no longer fall directly under
Proposition 4.1. In the case of elliptic curves ($g=1$), however, they could be
reduced to it by the considerations that follow.

One may of course vary Proposition 4.1 ad libitum by considering systems of
schemes over $k$ equipped with various structures. Suppose, for example, that
$X_0$ is an abelian scheme over $k$ with a marked origin (i.e. $X_0$ is
regarded as a group scheme over $k$), and let $F(A)$ be the set of isomorphism
classes of abelian schemes over $A$ (i.e. group schemes proper and smooth over
$A$) equipped with an isomorphism of abelian schemes $(*)$. One verifies that
imposing a multiplicative structure (or even only a “unit section”) eliminates
infinitesimal automorphisms, and consequently that a formal moduli scheme
exists, corresponding to a complete Noetherian local ring $O$. Moreover, if
$X$ is a proper smooth scheme with “absolutely connected” fibers over a locally
Noetherian prescheme $S$, every multiplicative structure on $X$ having a unit
section is necessarily associative and commutative (provided at least that it
is so on one fiber and that $S$ is connected), and it is uniquely determined
by the knowledge

\src{387}
of the unit section. Furthermore, suppose that $S$ is the spectrum of an
Artinian local ring $A$ with residue field $k$, that $X$ is proper over $A$ and
equipped with a section $s$, and that $X\otimes_Ak$ is equipped with the
structure of an abelian scheme over $k$ whose identity element is the point of
$X\otimes_Ak$ corresponding to $s$. A straightforward obstruction calculation,
together with an argument due to Serre, proves that $X$ does in fact possess a
multiplicative structure having $s$ as its unit section. (Starting from this,
one can use the “existence theorem” of [2] to pass to the case where $A$ is
complete Noetherian local, and then the descent technique of [3] for the
general case, to prove the analogous statement for every connected locally
Noetherian $S$.) This proves that the functor $F(A)$ considered here is
isomorphic to the analogous functor defined at the beginning of this number by
forgetting the multiplicative structure on $X_0$. In particular, if
$\mathfrak m$ is the maximal ideal of $O$, then $\mathfrak m/\mathfrak m^2$ is
canonically isomorphic to the dual of $H^1(X_0,\Theta_{X_0/k})$ and hence has
dimension $n^2$, where $n=\dim X_0$. It would be very interesting to determine
whether $O$ is smooth over $\Lambda$, i.e. isomorphic to a formal power-series
algebra in $n^2$ indeterminates over $\Lambda$. Paragraph 1, Proposition 5.2,
gives an equivalent formulation of this question as an existence problem for
abelian schemes reducing to a given abelian scheme. Transcendental methods show
in any case that the answer is affirmative if $k$ has characteristic zero. In
characteristic $p\ne0$, one clearly need only consider the case in which
$\Lambda$ is the ring of Witt vectors constructed on an algebraically closed
field $k$. Perhaps the time has come for the “Greenberg functor” to prove its
worth\ldots\footnote{The 1962 erratum reports that the formal moduli scheme of
an abelian variety over a field is smooth over the Witt ring: every abelian
scheme over an Artinian local quotient lifts to the ring of which it is a
quotient. The proof uses the variance properties of the lifting obstruction
class introduced in Exposé 182, p. 12. It also recalls Mumford's constructions
of moduli schemes for curves of genus $g$ and for polarized abelian schemes
(cf. SHMT).}

\subsection*{5. Extension of coverings}
\addcontentsline{toc}{subsection}{5. Extension of coverings}

Let $\mathfrak X$ be a Noetherian formal prescheme [2], and let $U$ be an open
subset of $\mathfrak X$, locally defined by the “nonvanishing” of a section of
$\mathcal O_{\mathfrak X}$ that is a non-zero-divisor not only on
$\mathfrak X$ but on every $X_n$,\footnote{This qualification is imposed by
the 1962 erratum.} hence large enough that every section of
$\mathcal O_{\mathfrak X}$ on an open subset $V$ which vanishes on $U\cap V$
is zero. Let $\mathcal J$ be an “ideal of definition” for $\mathfrak X$, and
put
\[
  X_n=(\mathfrak X,\mathcal O_{\mathfrak X}/\mathcal J^{n+1}),
\]
an ordinary Noetherian prescheme. If $\mathfrak X'$ and $\mathfrak X''$ are
two flat coverings of $\mathfrak X$ (i.e. two preschemes over $\mathfrak X$
defined by sheaves of algebras that are coherent and locally free as sheaves of
modules), unramified over $U$, the evident map
\[
  \operatorname{Hom}_{\mathfrak X}(\mathfrak X',\mathfrak X'')
  \longrightarrow
  \operatorname{Hom}_{X_0}(X'_0,X''_0)
\]
is injective. In particular, an automorphism of $\mathfrak X'$ inducing the
identity on $X'_0$ is the identity. This permits the descent technique to be
applied to the situation. In particular, start with a flat covering $X'_0$ of
$X_0$, unramified

\src{388}
over $U_0$, and let $G(\mathfrak X)$ be the set of isomorphism classes (with
the isomorphisms inducing the identity on $X'_0$) of flat coverings
$\mathfrak X'$ of $\mathfrak X$ that induce $X'_0$ on $X_0$ (and hence are
necessarily unramified over $U$). Define $G(V)$ similarly for every open subset
$V$ of $\mathfrak X$, and more generally define $G(\mathfrak Y)$ for every
formal prescheme $\mathfrak Y$ over $\mathfrak X$. The results of [2] and [3]
first imply the following:

\noindent a. As $V$ ranges over the open subsets of $\mathfrak X$, the $G(V)$
form a sheaf on $\mathfrak X$, denoted $G_{\mathfrak X}=G$. The restriction of
this sheaf to $U$ is the constant sheaf whose fibers each consist of one
element.

More generally, determining the fibers of $G_{\mathfrak X}$ is a question
about complete local rings. More precisely:

\noindent b. For every $x\in\mathfrak X$ one has
\[
  G_x=G\bigl(\operatorname{Spec}(\mathcal O_{\mathfrak X,x})\bigr)
  \subset
  G\bigl(\operatorname{Spec}(\widehat{\mathcal O}_{\mathfrak X,x})\bigr),
\]
(the latter consists of isomorphism classes of finite free algebras $B$ over
$\widehat{\mathcal O}_{\mathfrak X,x}$, equipped with an isomorphism of
$B\otimes_{\widehat{\mathcal O}_{\mathfrak X,x}}\mathcal O_{X_0,x}$ with
$\widehat{\mathcal O}'_{0x}$, where $\mathcal O'_0$ is the sheaf of algebras on
$X_0$ that defines $X'_0$).\footnote{The printed formula identifies $G_x$
directly with the completed term and omits the hat on
$\mathcal O'_{0x}$. The inclusion and completion are prescribed by the 1962
erratum.}

\noindent c. There is a canonical isomorphism
\[
  G_{\mathfrak X}=\varprojlim_nG_{X_n};
\]
in other words, for every open subset $V$ of $\mathfrak X$,
\[
  G(V)=\varprojlim_nG(V_n).
\]

\noindent d. Suppose that $\mathfrak X$ is obtained from an ordinary proper
scheme $X$ over a complete Noetherian local ring $\Lambda$ with an ideal of
definition $\mathfrak m$, by taking the $\mathcal J$-adic completion of
$\mathcal O_X$, where $\mathcal J=\mathfrak m\mathcal O_X$. Then
$G(\mathfrak X)$ is canonically isomorphic to the set of classes of flat
coverings of the ordinary scheme $X$ that “reduce to $X'_0$.”

Figuratively, (a) and (b) establish the fundamental relations between the
local and global aspects of the problem, (c) gives the relations between the
“finite” and “infinitesimal” aspects, and finally (d) recalls, under the stated
conditions, the identity between the “formal” and “algebraic” aspects.

Now suppose that $\mathfrak X$ is defined over a complete Noetherian local ring
$\Lambda$, with $\mathcal J=\mathfrak m\mathcal O_{\mathfrak X}$ (so that
$X_0$ is a prescheme over $\Lambda/\mathfrak m$). For every algebra $A$ finite
over $\Lambda$, put
\[
  F(A)=G(\mathfrak X\times_\Lambda A),
\]
where the marking is taken after base change to the reduced coefficient algebra $A_{\mathrm{red}}$, not literally to $A/\mathfrak mA$. This is a covariant functor of $A$ with values in the category of sets, and by
(c) it is completely known if it is known for Artinian $A$. Saying that this
functor is pro-representable is equivalent to saying that it has the form
\[
  F(A)=\operatorname{Hom}_{\text{topological }\Lambda\text{-algebras}}(O,A),
\]

\src{389}
where $O$ is a topological $\Lambda$-algebra of the type considered in A,
paragraph 5; it is also equivalent to saying the same after restricting to
Artinian $A$. The conjunction of Theorems 1 and 2 therefore gives:

\result{PROPOSITION}{5.1} Provided that every Artinian base-changed family classified by the functor is flat over its coefficient algebra and that the chosen open remains schematically dense after the relevant base changes, the preceding functor is pro-representable.

Of course, by (a), if $U=\mathfrak X$, then $G(\mathfrak Y)$ consists of a
single element for every $\mathfrak Y$ over $\mathfrak X$, and the functor $F$
is then of little interest (one has $O=\Lambda$). In practically every other
case, it appears that the topological local ring $O$ is not Noetherian. Its
existence nevertheless strikingly exhibits the “continuous” nature of the set
$G(\mathfrak X)$ of solutions (corresponding intuitively to the existence of a
“continuous” choice in the way ramification propagates when $X'_0$ is
extended). Compare this result with J.-P. Serre's viewpoint [4] in local class
field theory, which likewise brings out the continuous character of the
topological Galois group of the maximal abelian extension of a “geometric”
local field. Its dual group (in the sense of Pontryagin) appears there as an
inductive limit of algebraic, or at least quasi-algebraic, groups; here too,
extensions are classified by infinite-dimensional “varieties.” Moreover, in
the preceding discussion one may take $\mathfrak X$ to be the formal spectrum
of a complete local ring (of which $\Lambda$ will be, for example, a Cohen
subring), and one may hope that the results of this number can be used to study
extensions of a complete local ring of dimension $>1$. In both the local and
global cases they may make it possible to formulate precise relations between
higher ramification phenomena and phenomena in characteristic zero (which are
accessible by transcendental methods). In any case, it is the analysis
preceding Proposition 5.1 that makes it possible to extend the methods of [2]
for studying the fundamental group to the “tamely ramified” case and to solve
the “three-point problem” by transcendental methods.

Finally, note that the situation simplifies when $X_0$ has dimension one: by (a), together with formal gluing for the punctured completed local schemes,
\[
  G(\mathfrak X)\simeq
  \prod_iG\bigl(\operatorname{Spec}(\widehat{\mathcal O}_{\mathfrak X,x_i})\bigr),
\]
where the $x_i$ are the points of $X_0-U$: the “local” extensions at the
ramification points may be chosen arbitrarily. Normality of $X_0$ alone does not imply that the formal moduli scheme in 5.1 is smooth over $\operatorname{Spec}(\Lambda)$; the printed final assertion is false in general.

\src{390}
\begin{center}
  \textbf{BIBLIOGRAPHY}
\end{center}

\begin{description}
  \item[{[1]}] Grothendieck (Alexander). --- \emph{On some points of
    homological algebra}, Tohoku Math. J., vol. 9, 1957, pp. 119--221.
  \item[{[2]}] Grothendieck (Alexander). --- \emph{Formal geometry and
    algebraic geometry}, Bourbaki Seminar, vol. 11, 1958/59, fasc. 3, no. 182.
  \item[{[3]}] Grothendieck (Alexander). --- \emph{Descent technique and
    existence theorems in algebraic geometry, I: Generalities. Descent along
    faithfully flat morphisms}, Bourbaki Seminar, vol. 12, 1959/60, fasc. 1,
    no. 190.
  \item[{[4]}] Serre (Jean-Pierre). --- \emph{Local fields and isogenies},
    Bourbaki Seminar, vol. 11, 1958/59, fasc. 3, no. 185.
\end{description}

\par\smallskip
\begingroup\footnotesize
\endgroup
